\documentclass[11pt]{article}
\usepackage[margin=1in]{geometry}
\usepackage{amsmath,amssymb,amsthm}
\usepackage[colorlinks=true,linkcolor=blue,citecolor=blue,urlcolor=blue]{hyperref}
\newtheorem{theorem}{Theorem}[section]
\newtheorem{lemma}[theorem]{Lemma}
\newtheorem{corollary}[theorem]{Corollary}
\newtheorem{proposition}[theorem]{Proposition}
\newtheorem{remark}[theorem]{Remark}
\numberwithin{equation}{section}
\newcommand{\R}{\mathbb R}
\newcommand{\C}{\mathbb C}
\newcommand{\tr}{\operatorname{Tr}}
\newcommand{\vecop}{\operatorname{vec}}
\newcommand{\diag}{\operatorname{diag}}
\newcommand{\supp}{\operatorname{supp}}
\newcommand{\cR}{\mathcal R}
\title{Positive rational approximation of power divergences\texorpdfstring{\\}{ }and weighted negative powers}
\author{Working research manuscript}
\date{November 2026}
\begin{document}
\maketitle

\begin{abstract}
For $1\le a<b\le2$, we construct positive rational sums approximating
the power divergence $f_a$ uniformly relative to $f_b$ with error
$C_{a,b}\exp[-2\pi\sqrt{a(b-a)m/b}]$ using $m$ poles. We also construct
positive Stieltjes lower approximations to $x^{-c}$ on $(0,1]$ with
weighted error $C_{c,d}\exp[-2\pi\sqrt{(d-c)m}]$ for $0<c<d\le1$.
These are the complete bounds in Conjectures~1 and~2 of
Faust--Fawzi (2023), with fixed multiplicative constants at the
displayed exponents. The proofs combine the classical modified
Ganelius product estimate with a positive multipoint Gaussian rule
and a polynomial comparison used twice in its error integral.
We give the positive constructions, their exact rational degrees,
an arithmetic preparation procedure with root extraction charged
separately, and explicit semidefinite representations.
\end{abstract}

\noindent\textit{Review and priority status.}
This is a working manuscript. Its proof note and the complete typeset
source passed separate internal mathematical reviews, including an
exact-source audit. This is neither formal verification nor external
peer review. A bounded primary-source
literature check in November~2026 found no later resolution of the two
conjectures; it does not certify priority or exhaustive novelty.

\section{Statements and relation to earlier bounds}

For $\alpha\in[-1,2]\setminus\{0,1\}$ set
\begin{equation}\label{eq:falpha}
 f_\alpha(x)=\frac{x^\alpha-\alpha(x-1)-1}{\alpha(\alpha-1)},
 \qquad x>0.
\end{equation}
Use the continuous limits
$f_0(x)=x-1-\log x$ and $f_1(x)=x\log x-x+1$.
In particular, $f_\alpha(1)=f_\alpha'(1)=0$,
$f_\alpha''(x)=x^{\alpha-2}$ and $f_2(x)=(x-1)^2/2$.
The permitted positive rational approximations are
\begin{equation}\label{eq:allowed}
 \widehat f_m(x)=\sum_{i=1}^m u_i
       \frac{(x-1)^2}{1+T_i(x-1)},\qquad u_i\ge0,\quad T_i\in[0,1].
\end{equation}
Nodes may be sorted and taken distinct. Following
Faust--Fawzi~\cite{FF}, write $\epsilon^{[m]}_{a,b}$ for the infimum of
$\sup_{x>0}|f_a(x)-\widehat f_m(x)|/f_b(x)$ over these sums. At $x=1$
the ratio is interpreted by continuous extension.

\begin{theorem}[Power divergences]\label{thm:C1}
For each $1\le a<b\le2$ there is a finite constant $C_{a,b}$ such that
\begin{equation}\label{eq:C1}
 \epsilon^{[m]}_{a,b}\le C_{a,b}
       \exp\left[-2\pi\sqrt{\frac{a(b-a)}b\,m}\right]
       \qquad(m\ge1).
\end{equation}
Approximations attaining this bound up to the stated constant can be
chosen with $m$ distinct nodes in $(0,1)$ and strictly positive weights.
More explicitly, for every $L>0$ our construction has error at most
$C_{a,b}e^{-L}$ and uses at most
\begin{equation}\label{eq:C1count}
 \frac{bL^2}{4\pi^2a(b-a)}+3
\end{equation}
nodes. For $-1\le\beta<\alpha\le0$ it also gives
\begin{equation}\label{eq:C1negative}
 \epsilon^{[m]}_{\alpha,\beta}\le C_{1-\alpha,1-\beta}
  \exp\left[-2\pi\sqrt{\frac{(1-\alpha)(\alpha-\beta)}{1-\beta}\,m}\right].
\end{equation}
\end{theorem}

Let $\cR_{p,q}$ denote rational functions with numerator degree at most
$p$ and denominator degree at most $q$.

\begin{theorem}[Weighted negative powers]\label{thm:C2}
For each $0<c<d\le1$ there is a finite constant $C_{c,d}$ such that,
for every $m\ge1$, a rational function
\begin{equation}\label{eq:positiveStieltjes}
 R_m(x)=\sum_{i=1}^m\frac{w_i}{x+t_i},\qquad w_i>0,\quad t_i>0
\end{equation}
with distinct $t_i$ satisfies $0<R_m(x)\le x^{-c}$ for all $x>0$ and
\begin{equation}\label{eq:C2}
 \sup_{0<x\le1}x^d\bigl(x^{-c}-R_m(x)\bigr)
       \le C_{c,d}e^{-2\pi\sqrt{(d-c)m}}.
\end{equation}
Thus $R_m\in\cR_{m-1,m}$. For a target exponent $L>0$, the construction
has error at most $C_{c,d}e^{-L}$ with
\begin{equation}\label{eq:C2count}
 m\le\frac{L^2}{4\pi^2(d-c)}+\frac{L}{4\log2}+2.
\end{equation}
\end{theorem}

Theorem~\ref{thm:C1} proves Conjecture~1 of~\cite{FF}, including its
endpoint parameters. Setting $\alpha=-c$ and $\beta=-d$ in
Theorem~\ref{thm:C2} proves their Conjecture~2, with a stronger degree
restriction and positive lower approximants. All constants above may
depend on the fixed parameters. We claim neither matching lower bounds
nor an exact asymptotic multiplicative constant.

For comparison, \cite[Theorem~4]{FF} proves
$\epsilon^{[m]}_{a,b}\le C_{a,b}m^{3/2}e^{-\pi\sqrt{2a(b-a)m/b}}$.
Their Theorem~3 already gives the corresponding classical exponent
$2\pi\sqrt{\alpha(1-\alpha)m}$ in the different range $0<\alpha<1$.
Root-exponential approximation of positive powers and Markov functions
has a long history; see Stahl~\cite{Stahl} and Pekarskii~\cite{Pekarskii}.
Here the issue is the stated global weighted bounds with positive
residues and a fixed prefactor. The ingredients in the next section are
classical; their particular weighted use will give both results.

\section{A product estimate and finite-moment Gaussian quadrature}

We use the following form of the modified Ganelius product theorem.

\begin{lemma}[Classical product estimate]\label{lem:Ganelius}
For each fixed $0<r<1$ there is $A_r<\infty$ and, for each $n\ge1$,
explicit nodes $a_1,\ldots,a_n\in(0,1)$ such that
\begin{equation}\label{eq:Ganelius}
 \sup_{0\le s\le1}s^r
  \left|\prod_{j=1}^n\frac{s-a_j}{s+a_j}\right|
       \le A_r e^{-\pi\sqrt{rn}}.
\end{equation}
\end{lemma}

This is \cite[Definition~2.5 and Theorem~2.6]{TOS}, attributed there to
\cite[Lemma~1]{JH}. For clarity, its explicit nodes for sufficiently
large $n$ are as follows. Put
\[
 n_0=n-\left\lceil\frac\pi4\sqrt{nr}\right\rceil,
 \qquad \phi(v)=\exp\bigl(\pi\sqrt{v/r}\bigr).
\]
When $n_0>3$ and $n-n_0-1>0$, take
\begin{equation}\label{eq:Gnodes}
 a_j=\begin{cases}
 \phi(j-1)/\phi(n_0),&1\le j\le n_0,\\
 \phi(n_0-\tfrac12)/\phi(n_0),&j=n_0+1,\\
 1-\dfrac{j-n_0-1}{5(n-n_0-1)},&n_0+2\le j\le n.
 \end{cases}
\end{equation}
For the finitely many remaining $n$, choose $a_j=j/(n+1)$ and enlarge
$A_r$. Every use below has $0<r<1/2$. No uniform estimate in $r$ is
required. Coincident nodes, if any, are interpreted with multiplicity.

The measure used later need not have a finite moment of order $2m$.
We therefore include the precise quadrature argument.

\begin{lemma}[Gaussian rule with $2m-1$ finite moments]\label{lem:Gauss}
Let $\nu$ be a positive Borel measure on $(0,\infty)$ with infinite
support and finite moments through degree $2m-1$. There are distinct
$t_1,\ldots,t_m>0$ and $A_1,\ldots,A_m>0$ such that
\begin{equation}\label{eq:Gauss}
 \int P(t)\,d\nu(t)=\sum_{i=1}^m A_iP(t_i)
       \qquad(\deg P\le2m-1).
\end{equation}
Let $q_m$ be the monic degree-$m$ polynomial orthogonal to degrees
$0,\ldots,m-1$ for $\nu$. For $x>0$,
\begin{equation}\label{eq:CauchyError}
 \int\frac{d\nu(t)}{x+t}-\sum_i\frac{A_i}{x+t_i}
 =\frac1{q_m(-x)^2}\int\frac{q_m(t)^2}{x+t}\,d\nu(t).
\end{equation}
Moreover $q_m(t)/q_m(-x)$ minimizes
$\int p(t)^2(x+t)^{-1}\,d\nu(t)$ over real polynomials of degree at
most $m$ satisfying $p(-x)=1$.
\end{lemma}

\begin{proof}
The moment matrix for degrees $0,\ldots,m-1$ is positive definite,
so the orthogonality equations uniquely determine $q_m$. Suppose it
has fewer than $m$ sign-changing roots in $(0,\infty)$. Multiply it by
the product of those roots' linear factors. The resulting polynomial
has degree at most $2m-1$ and one nonzero sign on $(0,\infty)$.
Its finite integral is nonzero, contradicting orthogonality. Thus
$q_m$ has exactly $m$ distinct positive roots $t_i$.

Let $\ell_i$ be the corresponding Lagrange polynomials and set
$A_i=\int\ell_i\,d\nu$. Dividing a polynomial of degree at most
$2m-1$ by $q_m$ proves~\eqref{eq:Gauss}. Applying exactness to
$\ell_i^2$ gives $A_i=\int\ell_i^2\,d\nu>0$.

Apply~\eqref{eq:Gauss} to the polynomial
$[q_m(t)^2-q_m(-x)^2]/(t+x)$ to obtain~\eqref{eq:CauchyError}.
All its integrals are finite: the rational integrand on the right
grows only as $t^{2m-1}$. For a competitor $p$, write
\[
 p(t)-q_m(t)/q_m(-x)=(t+x)s(t),\qquad \deg s\le m-1.
\]
The cross term in the squared integral vanishes by orthogonality,
and the remaining added term is $\int(t+x)s(t)^2\,d\nu(t)\ge0$.
This proves the minimum assertion using only the stated moments.
\end{proof}

\begin{corollary}[Positive multipoint comparison]\label{cor:comparison}
Let $\mu$ be a positive measure on $(0,\infty)$ such that
$S_\mu(x)=\int(x+t)^{-1}\,d\mu(t)$ is finite for $x>0$.
Let $\Omega$ be a polynomial of degree at most $2m$, positive on
$(0,\infty)$, such that $d\nu=d\mu/\Omega$ satisfies
Lemma~\ref{lem:Gauss}. Define
\begin{equation}\label{eq:rescaledGauss}
 Q_m(x)=\sum_i\frac{A_i\Omega(t_i)}{x+t_i}.
\end{equation}
Then
\begin{equation}\label{eq:multipoint}
 S_\mu(x)-Q_m(x)=\frac{\Omega(-x)}{q_m(-x)^2}
               \int\frac{q_m(t)^2}{x+t}\,d\nu(t).
\end{equation}
For any real polynomial $q_0$ of degree at most $m$ with $q_0(-x)\ne0$,
\begin{equation}\label{eq:comparison}
 |S_\mu(x)-Q_m(x)|\le\frac{|\Omega(-x)|}{q_0(-x)^2}
             \int\frac{q_0(t)^2}{\Omega(t)(x+t)}\,d\mu(t).
\end{equation}
\end{corollary}

\begin{proof}
The quotient $[\Omega(t)-\Omega(-x)]/(t+x)$ has degree at most
$2m-1$, so Gaussian exactness reduces the error to $\Omega(-x)$
times~\eqref{eq:CauchyError}. Use the minimum assertion with
$p=q_0/q_0(-x)$ for~\eqref{eq:comparison}.
\end{proof}

\section{Power divergences: one global positive quadrature}

Put $H_a(x)=f_a(x)/(x-1)^2$, with $H_a(1)=1/2$, and
\begin{equation}\label{eq:Ka}
 K_a=\frac{\sin(\pi(a-1))}{\pi a(a-1)}\quad(1<a<2),
 \qquad K_1=1.
\end{equation}
For $1\le a<2$ we have
\begin{equation}\label{eq:Hrepresentation}
 H_a(x)=K_a\int_0^\infty\frac{t^a}{(1+t)^2(x+t)}\,dt.
\end{equation}
Indeed, the second derivative of $(x-1)^2/(x+t)$ is
$2(1+t)^2/(x+t)^3$. The beta integral gives
$2K_a\int_0^\infty t^a(x+t)^{-3}\,dt=x^{a-2}$, including $a=1$.
The value and first derivative at $1$ then identify $f_a$.
All differentiations are justified locally by integrable tails.
Also Taylor's integral formula implies, for $1\le b\le2$,
\begin{equation}\label{eq:weight}
 H_b(x)=\int_0^1(1-u)[1+u(x-1)]^{b-2}\,du
       \ge\tfrac12\min(1,x^{b-2}).
\end{equation}
In particular, this includes $b=2$ exactly.

Fix $1\le a<b\le2$, and write
\[
 \gamma=b-a,\qquad \kappa=\frac{a\gamma}{b},\qquad
 p=2-\frac{a+b}{2}\in(0,1).
\]
Take $n_-$ Ganelius nodes $a_j$ with parameter $a/4$, and $n_+$
Ganelius nodes $b_j$ with parameter $\gamma/4$. Define
\begin{align}\label{eq:Bcombined}
 B_-(s)&=\prod_{j=1}^{n_-}\frac{s-a_j}{s+a_j},&
 B_+(s)&=\prod_{j=1}^{n_+}\frac{s-b_j}{s+b_j},&
 B(s)&=B_-(s)B_+(1/s),\\
 D_-&=A_{a/4}^2e^{-\pi\sqrt{a n_-}},&
 D_+&=A_{\gamma/4}^2e^{-\pi\sqrt{\gamma n_+}}.\nonumber
\end{align}
Each product has modulus at most one on the positive axis. Squaring
\eqref{eq:Ganelius} therefore gives
\begin{equation}\label{eq:envelopes}
 B(s)^2\le
 \begin{cases}D_-s^{-a/2},&0<s\le1,\\
 D_+s^{\gamma/2},&s\ge1.
 \end{cases}
\end{equation}

Set $m=n_-+n_++1$ and
\begin{equation}\label{eq:OmegaC1}
 \Omega(t)=t(t+1)\prod_{j=1}^{n_-}(t+a_j)^2
                         \prod_{j=1}^{n_+}(t+1/b_j)^2,
 \qquad d\nu(t)=\frac{t^a\,dt}{(1+t)^2\Omega(t)}.
\end{equation}
This polynomial has degree $2m$. The density of $\nu$ is comparable
to $t^{a-1}$ near zero and to $t^{a-2m-2}$ near infinity.
Its moment of degree $2m-1$ thus has tail $t^{a-3}$ and is finite.
Apply Lemma~\ref{lem:Gauss} and set
\begin{equation}\label{eq:RC1}
 R_m(x)=K_a\sum_{i=1}^m\frac{A_i\Omega(t_i)}{x+t_i}.
\end{equation}
All residues are strictly positive, and all poles are distinct and negative.

In~\eqref{eq:comparison} choose the degree-$m$ polynomial
\[
 q_0(t)=(t-1)\prod_{j=1}^{n_-}(t-a_j)
                     \prod_{j=1}^{n_+}(t-1/b_j).
\]
Cancellation of its factors gives the global estimate
\begin{equation}\label{eq:GLOBAL}
 \frac{|H_a(x)-R_m(x)|}{K_a}
 \le\frac{x|x-1|}{(x+1)^2}B(x)^2
   \int_0^\infty B(t)^2
      \frac{t^{a-1}(t-1)^2}{(1+t)^3(x+t)}\,dt.
\end{equation}
The exact error has sign $\operatorname{sign}(x-1)$, apart from its
interpolation zeros, since $\Omega(-x)$ has that sign.
The two occurrences of $B^2$ in~\eqref{eq:GLOBAL} are both retained.

For $x\le1$, bound the outside factor by $x$ and split the integral
at $t=1$. Using~\eqref{eq:envelopes} yields
\begin{align*}
 \frac{|H_a(x)-R_m(x)|}{K_a}
 &\le D_-^2x^{1-a/2}\int_0^1\frac{t^{a/2-1}}{x+t}\,dt\\
 &\quad+D_-D_+x^{1-a/2}
                    \int_1^\infty\frac{t^{a-2+\gamma/2}}{x+t}\,dt\\
 &\le \frac{\pi D_-^2}{\sin(\pi a/2)}+\frac{D_-D_+}{p}.
\end{align*}
Here $0<a/2<1$ justifies extending the first beta integral to
$(0,\infty)$. In the second integral use $x+t\ge t$,
$x^{1-a/2}\le1$, and $p=2-a-\gamma/2>0$.
Dividing by~\eqref{eq:weight} proves
\begin{equation}\label{eq:LEFT}
 \frac{|H_a(x)-R_m(x)|}{H_b(x)}
 \le2K_a\left[\frac{\pi D_-^2}{\sin(\pi a/2)}
                         +\frac{D_-D_+}{p}\right],\quad x\le1.
\end{equation}

For $x\ge1$, the outside factor in~\eqref{eq:GLOBAL} is at most one.
The same split gives
\[
 \frac{|H_a(x)-R_m(x)|}{K_a}
 \le D_-D_+x^{\gamma/2}\frac{2}{ax}
   +D_+^2x^{\gamma/2}\int_1^\infty\frac{t^{-p}}{x+t}\,dt.
\]
The last integral is at most $\pi x^{-p}/\sin(\pi p)$.
After division by $H_b(x)\ge x^{b-2}/2$, the powers of $x$ in the
two terms are respectively $1-(a+b)/2<0$ and
$\gamma/2-p+2-b=0$. Hence
\begin{equation}\label{eq:RIGHT}
 \frac{|H_a(x)-R_m(x)|}{H_b(x)}
 \le2K_a\left[\frac2aD_-D_+
                         +\frac{\pi D_+^2}{\sin(\pi p)}\right],\quad x\ge1.
\end{equation}
All integrals remain finite at $a=1,b=2$.

\begin{proof}[Completion of Theorem~\ref{thm:C1}]
Equations~\eqref{eq:LEFT}--\eqref{eq:RIGHT} give a fixed-constant bound
by $D_-^2+D_-D_++D_+^2$. Set
\begin{equation}\label{eq:allocationL}
 n_-=\left\lceil\frac{L^2}{4\pi^2a}\right\rceil,
 \qquad n_+=\left\lceil\frac{L^2}{4\pi^2\gamma}\right\rceil.
\end{equation}
Each of these three terms is at most a fixed multiple of $e^{-L}$,
and $m=n_-+n_++1$ gives~\eqref{eq:C1count}.

For a prescribed sufficiently large integer $m$, instead take
\[
 n_-=\left\lfloor\frac\gamma b(m-1)\right\rfloor,
 \qquad n_+=m-1-n_-.
\]
Then $a n_- =\kappa m+O_{a,b}(1)$ and
$\gamma n_+=\kappa m+O_{a,b}(1)$. Thus each error term is bounded by
a fixed multiple of $e^{-2\pi\sqrt{\kappa m}}$.
For the finitely many smaller $m$, the same positive quadrature
construction with $\Omega(t)=t(t+1)^{2m-1}$ has finite relative error:
$H_a/H_b$ is bounded, and its finite positive Stieltjes sum divided by
$H_b$ is bounded at both endpoints and on compact intervals.
Enlarging $C_{a,b}$ covers these cases while retaining strict positivity.

Finally, define
\[
 T_i=\frac1{1+t_i},\qquad
 u_i=\frac{K_a A_i\Omega(t_i)}{1+t_i}>0.
\]
Then $(x-1)^2R_m(x)$ is exactly~\eqref{eq:allowed}, so the relative
error is unchanged. The double zero at $1$ is automatic. Gaussian
exactness on $\Omega(t)/(t+1)$ also gives $R_m(1)=H_a(1)=1/2$.

For the negative-parameter assertion use
$f_\alpha(x)=x f_{1-\alpha}(1/x)$. The corresponding transformation
of a kernel in~\eqref{eq:allowed} replaces $T_i$ by $1-T_i$ and
leaves its weight unchanged. Substitute $x\mapsto1/x$ in the relative
error and apply~\eqref{eq:C1} with $a=1-\alpha$, $b=1-\beta$.
\end{proof}

\section{Weighted negative powers: a positive lower approximation}

\begin{proof}[Proof of Theorem~\ref{thm:C2}]
Put $\gamma=d-c$, $r=\gamma/4$, and $K_c=\sin(\pi c)/\pi>0$. The
Stieltjes identity is
\begin{equation}\label{eq:powerStieltjes}
 x^{-c}=K_c\int_0^\infty\frac{t^{-c}}{x+t}\,dt.
\end{equation}
Take $n$ nodes $a_j$ from Lemma~\ref{lem:Ganelius}, write
$B_n(s)=\prod_j(s-a_j)/(s+a_j)$, and choose $k\ge1$. Set $m=n+k$ and
\[
 \Omega(t)=(1+t)^{2k}\prod_{j=1}^n(t+a_j)^2,
 \qquad d\nu(t)=\frac{t^{-c}}{\Omega(t)}\,dt.
\]
The moment of degree $2m-1$ has tail $t^{-c-1}$, and the density
near zero is comparable to $t^{-c}$. All required moments are finite.
Rescaling its Gaussian rule as in~\eqref{eq:rescaledGauss}, including
the factor $K_c$, produces~\eqref{eq:positiveStieltjes}. Since
$\Omega(-x)\ge0$, \eqref{eq:multipoint} gives $R_m(x)\le x^{-c}$
for all $x>0$.

Use $q_0(t)=\prod_j(t-a_j)$, of degree $n\le m$, in the comparison.
For $0<x\le1$,
\begin{equation}\label{eq:C2integral}
 \frac{x^d(x^{-c}-R_m(x))}{K_c}
 \le x^d(1-x)^{2k}B_n(x)^2
      \int_0^\infty\frac{B_n(t)^2t^{-c}}{(1+t)^{2k}(x+t)}\,dt.
\end{equation}
Let $D=A_r^2e^{-\pi\sqrt{\gamma n}}$. On $(0,1]$ the product bound
gives $B_n(s)^2\le D s^{-\gamma/2}$.
For $t\le1$ use this bound twice and discard the remaining factors
at most one. Its contribution to~\eqref{eq:C2integral} is at most
\[
 D^2x^{d-\gamma/2}\int_0^1\frac{t^{-c-\gamma/2}}{x+t}\,dt
       \le\frac{\pi D^2}{\sin(\pi(c+d)/2)}.
\]
The beta integral is finite since $0<(c+d)/2<1$; its residual power
of $x$ is $d-c-\gamma=0$. For $t\ge1$, use $|B_n(t)|\le1$,
$(1+t)^{-2k}\le2^{-2k}$ and $(x+t)^{-1}\le t^{-1}$.
As $x^{d-\gamma/2}\le1$, the contribution is at most $D2^{-2k}/c$.
Consequently
\begin{equation}\label{eq:C2estimate}
 \sup_{0<x\le1}x^d(x^{-c}-R_m(x))
 \le K_c\left[\frac{\pi D^2}{\sin(\pi(c+d)/2)}
                                     +\frac Dc\,2^{-2k}\right].
\end{equation}
Choose
\[
 n=\left\lceil\frac{L^2}{4\pi^2\gamma}\right\rceil,
 \qquad k=\left\lceil\frac{L}{4\log2}\right\rceil.
\]
Both terms are bounded by a fixed multiple of $e^{-L}$, proving
\eqref{eq:C2count}.

For a prescribed large $m$, let $L>0$ solve
$L^2/(4\pi^2\gamma)+L/(4\log2)+2=m$.
The preceding choices use at most $m$ nodes; increase $k$ to use
the full budget if necessary. Inequality~\eqref{eq:C2estimate} only
improves under this increase. The positive root satisfies
$L=2\pi\sqrt{\gamma m}-O_{c,d}(1)$, which proves~\eqref{eq:C2} with
a fixed prefactor. For the finitely many remaining $m$, use the
Gaussian rule for $d\nu=t^{-c}(1+t)^{-2m}dt$ and rescale its weights.
It gives $m$ distinct positive poles, a strictly positive global lower
approximation, and weighted error at most $x^{d-c}\le1$ on $(0,1]$.
Increasing the fixed constant handles these cases.
\end{proof}

\section{Arithmetic preparation and an error certificate}

Both constructions can be prepared using elementary interpolation
data rather than integration of the modified moments. The following
cost statement uses exact real arithmetic; root extraction is a
separate operation.

\begin{proposition}\label{prop:preparation}
For the power-divergence construction with $m=n_-+n_++1$, $R_m$ is
the unique type $[m-1/m]$ rational Hermite interpolant to $H_a$ at
$0$ once, $1$ once, the $a_j$ twice, and the $1/b_j$ twice.
For the negative-power construction with $m=n+k$, $R_m$ is the
corresponding interpolant to $x^{-c}$ at each $a_j$ twice and at $1$
with multiplicity $2k$.

For either construction, computing the numerator and monic denominator
costs $O(m^3)$ arithmetic operations and $O(m^2)$ storage, plus
$O(m)$ elementary power, logarithm, exponential and square-root
evaluations. Recovering the positive-kernel coefficients additionally
requires the roots of a degree-$m$ polynomial; after that, residue
evaluation costs $O(m^2)$ arithmetic operations.
\end{proposition}

\begin{proof}
The zeros of $\Omega(-x)$ in~\eqref{eq:multipoint} give the positive
Hermite data, since its Cauchy-error factor is analytic there.
For the power-divergence construction, the boundary value at zero
follows directly from Gaussian exactness on $\Omega(t)/t$, a polynomial
of degree $2m-1$. The value at one follows from exactness on
$\Omega(t)/(t+1)$. Here $H_a(0)=1/a$ and $H_a(1)=1/2$.
All the other values and derivatives are obtained from $x^a$, or
$x\log x$ when $a=1$. Repeated Ganelius nodes mean combined Hermite
multiplicities; their higher derivatives can be generated algebraically
within the stated arithmetic bound.

Each positive sum has $m$ distinct poles and nonzero residues, so its
numerator and degree-$m$ monic denominator are coprime. The $2m$
interpolation conditions are linear in their $2m$ unknown coefficients.
They have a unique solution: cross-multiplication of two solutions
gives a polynomial of degree at most $2m-1$ with $2m$ zeros counted
with multiplicity. It is zero, and coprimality and normalization then
identify the coefficient pairs. This argument also excludes spurious
linear-system solutions whose denominator vanishes at a data point.
Dense elimination proves the coefficient cost. If $P/Q$ is the result,
the poles are the negative simple roots $-t_i$ of $Q$ and the residues
are $P(-t_i)/Q'(-t_i)$.
\end{proof}

The constants in Lemma~\ref{lem:Ganelius} need not be known numerically
to certify an individual power-divergence approximation. Replace $D_-$
and $D_+$ in~\eqref{eq:LEFT}--\eqref{eq:RIGHT} by the actual maxima
\[
 \max_{0\le s\le1}s^{a/2}B_-(s)^2,
 \qquad \max_{0\le s\le1}s^{\gamma/2}B_+(s)^2.
\]
These give the same rigorous error bounds. For nodes $z_j$, the positive
critical points of $s^{2r}\prod_j[(s-z_j)/(s+z_j)]^2$ satisfy
\begin{equation}\label{eq:certificatePolynomial}
 r\prod_j(s^2-z_j^2)+2s\sum_j z_j
                       \prod_{\ell\ne j}(s^2-z_\ell^2)=0.
\end{equation}
Thus each maximum is determined by endpoints and roots of a polynomial
of degree at most twice the number of nodes, excluding the product's
zeros when taking logarithmic derivatives. The analogous maximum
certifies~\eqref{eq:C2estimate}. These are additional polynomial root
problems, explicitly charged when this certificate is used.

The proposition is a coefficient-arithmetic statement plus the specified
root problems. It asserts no numerical stability, machine-precision
accuracy, or certified bit-complexity bound. Exact positive existence
and the approximation theorems do not depend on a particular numerical
implementation of root extraction.

\section{Explicit semidefinite consequences}

We give representations whose block sizes follow directly from Schur
complements. This avoids dependence on the smaller quantum block
representation stated without proof in~\cite[Theorem~5]{FF}.

For scalars $x>0$, the epigraph $\tau\ge\widehat f_m(x)$ is represented
by $m$ auxiliary scalars $y_i$ and $m$ blocks of size $2$:
\begin{equation}\label{eq:scalarSDP}
 \tau\ge\sum_i u_i y_i,\qquad
 \begin{pmatrix}1+T_i(x-1)&x-1\\x-1&y_i\end{pmatrix}\succeq0.
\end{equation}
For positive definite $N\times N$ Hermitian $A,B$, define the matrix
perspective
\[
 P_f(A,B)=B^{1/2}f(B^{-1/2}AB^{-1/2})B^{1/2}.
\]
Its rational approximation is
\[
 P_{\widehat f_m}(A,B)=\sum_i u_i(A-B)
                      [(1-T_i)B+T_iA]^{-1}(A-B).
\]
The epigraph in the Loewner order has the exact representation
\begin{equation}\label{eq:matrixSDP}
 E\succeq\sum_i u_iY_i,\qquad
 \begin{pmatrix}(1-T_i)B+T_iA&A-B\\ A-B&Y_i\end{pmatrix}\succeq0,
\end{equation}
with $m$ Hermitian auxiliary matrices and $m$ blocks of size $2N$.
Scalar functional calculus followed by congruence gives
\begin{equation}\label{eq:matrixError}
 -\epsilon P_{f_b}(A,B)\preceq P_{f_a}(A,B)-P_{\widehat f_m}(A,B)
                          \preceq\epsilon P_{f_b}(A,B)
\end{equation}
whenever the scalar relative error is at most $\epsilon$.

There is also an explicit representation for the Petz divergence
\[
 S_f(\rho\Vert\sigma)=
  \langle\sigma^{1/2},f(L_\rho R_\sigma^{-1})\sigma^{1/2}\rangle_{\rm HS},
\]
where $L_\rho(Z)=\rho Z$, $R_\sigma(Z)=Z\sigma$, and
$\rho,\sigma\succ0$. The two operators commute and are positive
self-adjoint on the Hilbert--Schmidt space.
For the kernel $f_T(x)=(x-1)^2/[1+T(x-1)]$, put
\[
 v=\vecop(\rho-\sigma),\qquad
 M_T=T I\otimes\rho+(1-T)\sigma^{\mathsf T}\otimes I.
\]
Diagonalizing the left and right multiplication operators gives
\begin{equation}\label{eq:PetzKernel}
 S_{f_T}(\rho\Vert\sigma)=v^*M_T^{-1}v.
\end{equation}
Indeed, its summands are
$|\langle u_i,v_j\rangle|^2(\rho_i-\sigma_j)^2/
[T\rho_i+(1-T)\sigma_j]$ in eigenbases of $\rho$ and $\sigma$.
Thus $\tau\ge S_{\widehat f_m}(\rho\Vert\sigma)$ is represented by
\begin{equation}\label{eq:PetzSDP}
 \tau\ge\sum_i u_i y_i,\qquad
 \begin{pmatrix}M_{T_i}&v\\v^*&y_i\end{pmatrix}\succeq0.
\end{equation}
There are exactly $m$ blocks of size $N^2+1$ and $m$ auxiliary scalars.
Every block is affine in $\rho,\sigma,y_i$, so this is a convex SDP
epigraph representation.

\begin{corollary}[Relative entropy approximation]\label{cor:entropy}
For trace-one positive definite $\rho,\sigma$, let
$D(\rho\Vert\sigma)=\tr[\rho(\log\rho-\log\sigma)]$.
There are convex functions $D_m=S_{\widehat f_m}$ with the
representation~\eqref{eq:PetzSDP} such that
\begin{equation}\label{eq:entropyError}
 |D(\rho\Vert\sigma)-D_m(\rho\Vert\sigma)|
 \le\frac{C e^{-\pi\sqrt{2m}}}{2}
                  \bigl(\tr[\rho^2\sigma^{-1}]-1\bigr).
\end{equation}
For scalar relative error $\epsilon$, the construction uses at most
$\log^2(C/\epsilon)/(2\pi^2)+3$ blocks for sufficiently small $\epsilon$.
\end{corollary}

\begin{proof}
Use Theorem~\ref{thm:C1} with $a=1,b=2$. Functional calculus for
$L_\rho R_\sigma^{-1}$ bounds the absolute divergence error by
$\epsilon S_{f_2}(\rho\Vert\sigma)$. For trace-one matrices,
$S_{f_1}=D$ and $S_{f_2}=(\tr[\rho^2\sigma^{-1}]-1)/2$.
The node count follows from~\eqref{eq:C1count} with
$L=\log(C/\epsilon)$.
\end{proof}

These statements concern approximation size and explicit representations.
They do not assert an SDP solver runtime, nor do they establish the
smaller $2N$ quantum block representation mentioned in~\cite{FF}.

\section*{Bibliographic scope of this working version}

The primary arXiv record of~\cite{FF} inspected in November~2026 lists
only its 2023 first version. Faust's December~2023 thesis still states
the two questions as Conjectures~3.1 and~3.2~\cite{FaustThesis}; the
current author publication page lists the joint paper as a
preprint~\cite{FawziPage}. Targeted searches for these conjectures and
the corresponding weighted negative-power problem did not locate a
later primary resolution. This is a bounded screen, and no claim of an
exhaustive priority determination is made. The proof uses the classical
product estimate with explicit attribution and proves the required
Gaussian comparison directly.

\begin{thebibliography}{9}\raggedright
\bibitem{FF}
O.~Faust and H.~Fawzi.
\emph{Rational approximations of operator monotone and operator convex functions}.
arXiv:2305.12405v1, 2023.
\url{https://arxiv.org/abs/2305.12405}.

\bibitem{TOS}
K.~Tanaka, T.~Okayama, and M.~Sugihara.
An optimal approximation formula for functions with singularities.
\emph{Journal of Approximation Theory} \textbf{234} (2018), 82--107.
\href{https://doi.org/10.1016/j.jat.2018.06.004}{doi:10.1016/j.jat.2018.06.004}.
The node and theorem numbering used here is from arXiv:1610.06844v1,
\url{https://arxiv.org/pdf/1610.06844}.

\bibitem{JH}
A.~P.~Jang and S.~Haber.
Numerical indefinite integration of functions with singularities.
\emph{Mathematics of Computation} \textbf{70} (2001), 205--221.
\href{https://doi.org/10.1090/S0025-5718-00-01226-6}{doi:10.1090/S0025-5718-00-01226-6}.

\bibitem{Pekarskii}
A.~A.~Pekarskii.
Best uniform rational approximations of Markov functions.
\emph{Algebra i Analiz} \textbf{7}(2) (1995), 121--132;
English translation, \emph{St.~Petersburg Mathematical Journal}
\textbf{7}(2) (1996), 277--285.
\url{https://www.mathnet.ru/eng/aa551}.

\bibitem{Stahl}
H.~R.~Stahl.
Best uniform rational approximation of $x^\alpha$ on $[0,1]$.
\emph{Acta Mathematica} \textbf{190} (2003), 241--306.
\href{https://doi.org/10.1007/BF02392691}{doi:10.1007/BF02392691}.

\bibitem{FaustThesis}
O.~P.~B.~Faust.
\emph{Convex Relaxations: Beyond Polynomials, Splitting Methods and Average Case Analysis}.
PhD thesis, University of Cambridge, December 2023, Chapter~3.
\href{https://api.repository.cam.ac.uk/server/api/core/bitstreams/95b542c8-984c-45ca-bfee-e90bcad55a73/content}{Cambridge repository full text}.

\bibitem{FawziPage}
H.~Fawzi. Publications. Author website, inspected November 2026.
\url{https://www.damtp.cam.ac.uk/user/hf323/publications.html}.
\end{thebibliography}
\end{document}
